\input{../tex-inputs/latex-leseansicht-vorspann}

\section[Rainer Maria Rilke an Arthur Schnitzler, 3. 12. 1901]{L04540 Rainer Maria Rilke an Arthur Schnitzler, 3. 12. 1901}
\nopagebreak\mylabel{L04540v}
\rehead{ }\normalsize\beginnumbering\briefempfaengerindex{Schnitzler, Arthur@\textsc{Schnitzler, Arthur}!zzzRilke, Rainer Maria@\emph{von Rainer Maria Rilke}!1901-12-031@{3. 12. 1901}|(be}
\toendnotes[C]{\smallbreak\pagebreak[2]}
\correspDesc{Versand  durch Rainer Maria Rilke am 3. 12. 1901 in Westerwede
\newline{}Erhalt  durch Arthur Schnitzler im Zeitraum [4. 12. 1901 – 8. 12. 1901?] in Wien}\toendnotes[C]{\smallbreak}
\Standort{CUL, Schnitzler, B 84.}
\physDesc{Brief, 1 Blatt, 4 Seiten, 3257 Zeichen
\newline{}Handschrift: schwarze Tinte, deutsche Kurrent
\newline{}Schnitzler: mit rotem Buntstift eine Unterstreichung }
\buchAbdrucke{\weitereDrucke{1) \emph{Rainer Maria Rilke und Arthur Schnitzler. Ihr Briefwechsel.}. Mit Anmerkungen herausgegeben von Heinrich Schnitzler In: \emph{Wort und Wahrheit}, Jg. 13, Nr. 4, April 1958, S. 282–298, hier S. 286–287.} \weitereDrucke{2) \emph{Rainer Maria Rilke und Arthur Schnitzler. Ihr Briefwechsel.}. Mit Anmerkungen herausgegeben von Heinrich Schnitzler In: \emph{Wort und Wahrheit}, Jg. 13, Nr. 4, April 1958, S. 282–298, hier S. 288–289.} \weitereDrucke{3) Rainer Maria Rilke: \emph{Briefe in zwei Bänden}. Herausgegeben von Horst Nalewski. Frankfurt am Main, Leipzig: \emph{Insel} 1991, S. XXXX.} \weitereDrucke{4) Rainer Maria Rilke: \emph{Briefe zur Politik}. Herausgegeben von Joachim W. Storck. Frankfurt am Main, Leipzig: \emph{Insel} 1992, S. XXXX.} }\toendnotes[C]{\smallbreak}
\pstart
           {\pb}\textsc{Westerwede bei Bremen\oindex{Westerwede@\textbf{Westerwede}|pw}},\pend
           
\pstart
           am 3. Dezember 1901.\pend
           
\pstart{}Lieber und verehrter Doctor \textsc{Schnitzler},\pend\vspace{0.5em}
\pstart
           an dieſem Abend, vor einer halben Stunde habe ich meiner Frau\pwindex{Westhoff, Clara 21.\,11.\,1878 Bremen – 9.\,3.\,1954 Fischerhude@\textsc{Westhoff, Clara} (21.\,11.\,1878 Bremen – 9.\,3.\,1954 Fischerhude)|pwv} das kleine \label{K_L04540-1v}\edtext{Schauſpiel »\textsc{Lebendige Stunden\pwindex{Schnitzler, Arthur 15.\,5.\,1862 Wien – 21.\,10.\,1931 ebd.@\textsc{Schnitzler, Arthur} (15.\,5.\,1862 Wien – 21.\,10.\,1931 ebd.), \emph{Schriftsteller, Mediziner}!Lebendige Stunden@\strich\emph{Lebendige Stunden}|pw}}«, das wir in der \textsc{Neuen Deutschen Rundschau\pwindex{Neue Deutsche Rundschau@\emph{Neue Deutsche Rundschau}|pw}}}{\lemma{\textnormal{\emph{Schauspiel … Rundschau}}}\Cendnote{\textnormal{Arthur Schnitzler: \emph{Lebendige
                           Stunden. Schauspiel in einem Aufzug}\pwindex{Schnitzler, Arthur 15.\,5.\,1862 Wien – 21.\,10.\,1931 ebd.@\textsc{Schnitzler, Arthur} (15.\,5.\,1862 Wien – 21.\,10.\,1931 ebd.), \emph{Schriftsteller, Mediziner}!Lebendige Stunden@\strich\emph{Lebendige Stunden}|pwk}. In: \emph{Neue Deutschen Rundschau}\pwindex{Neue Deutsche Rundschau@\emph{Neue Deutsche Rundschau}|pwk}, Jg. 12, H. 12,
                     Dezember 1901, S. 1297–1306.}}}\label{K_L04540-1}{ }\introOben{}fanden\introOben{}, vorgeleſen – und es hat uns beide{ }ſehr unmittelbar ergriffen und
               feſtgehalten und – hält uns noch. Mir iſt es ganz beſonders lieb um einer gewiſſen,
               man könnte faſt{ }ſagen: äußerlichen Verwandtſchaft mit »\textsc{Der Gefährtin\pwindex{Schnitzler, Arthur 15.\,5.\,1862 Wien – 21.\,10.\,1931 ebd.@\textsc{Schnitzler, Arthur} (15.\,5.\,1862 Wien – 21.\,10.\,1931 ebd.), \emph{Schriftsteller, Mediziner}!Gefährtin. Schauspiel in einem Akt@\strich\emph{Die Gefährtin. Schauspiel in einem Akt}|pw}}« willen. Und jetzt{ }ſcheint mir{ }ſogar, als ob dieſe Verwandtſchaft nicht{ }ſo
               äußerlich wäre, wie man zunächſt meinen könnte. Denn es iſt eine ganz beſtimmte
               Stimmung, gleichſam {\pb}der Duft vom Weſen
               einer toten Frau, in dem dieſe beiden Dramen\pwindex{Schnitzler, Arthur 15.\,5.\,1862 Wien – 21.\,10.\,1931 ebd.@\textsc{Schnitzler, Arthur} (15.\,5.\,1862 Wien – 21.\,10.\,1931 ebd.), \emph{Schriftsteller, Mediziner}!Lebendige Stunden@\strich\emph{Lebendige Stunden}|pwv}\pwindex{Schnitzler, Arthur 15.\,5.\,1862 Wien – 21.\,10.\,1931 ebd.@\textsc{Schnitzler, Arthur} (15.\,5.\,1862 Wien – 21.\,10.\,1931 ebd.), \emph{Schriftsteller, Mediziner}!Gefährtin. Schauspiel in einem Akt@\strich\emph{Die Gefährtin. Schauspiel in einem Akt}|pwv}{ }ſpielen wie
               andere in einem Abend oder in einer bürgerlichen Wohnſtube. Der leiſe Nachklang einer
               Perſönlichkeit iſt hier dramatiſiert und wirkt geradezu wie der Ort, an dem die Szene
               vor{ }ſich geht. Der leere Lehnſtuhl iſt der Platz an dem die Lebende oder Sterbende zu{ }ſuchen war: das Stück\pwindex{Schnitzler, Arthur 15.\,5.\,1862 Wien – 21.\,10.\,1931 ebd.@\textsc{Schnitzler, Arthur} (15.\,5.\,1862 Wien – 21.\,10.\,1931 ebd.), \emph{Schriftsteller, Mediziner}!Lebendige Stunden@\strich\emph{Lebendige Stunden}|pw} aber{ }ſpielt in jener Stunde (jener
               einen geheimnisvollen Stunde), wo die Tote gleichſam in allen Dingen iſt, wo alles
               erfüllt iſt von ihr, wie erſchüttert von Ihrem Abſchied. Und durch dieſe\strikeout{n} Mitergriffenheit des Raumes haben sie einen Akteur
               aus ihm gemacht, deſſen Schweigen die\strikeout{\textcolor{gray}{s}} Stimmen der Redenden wunderſam im Gleichgewichte hält. Und der Konflikt dieſes
               Dramas\pwindex{Schnitzler, Arthur 15.\,5.\,1862 Wien – 21.\,10.\,1931 ebd.@\textsc{Schnitzler, Arthur} (15.\,5.\,1862 Wien – 21.\,10.\,1931 ebd.), \emph{Schriftsteller, Mediziner}!Lebendige Stunden@\strich\emph{Lebendige Stunden}|pw} liegt noch tiefer als der der Gefährtin\pwindex{Schnitzler, Arthur 15.\,5.\,1862 Wien – 21.\,10.\,1931 ebd.@\textsc{Schnitzler, Arthur} (15.\,5.\,1862 Wien – 21.\,10.\,1931 ebd.), \emph{Schriftsteller, Mediziner}!Gefährtin. Schauspiel in einem Akt@\strich\emph{Die Gefährtin. Schauspiel in einem Akt}|pw} und iſt noch dramatiſcher, weil die letzten heftigen
               Bewegungen der Handlung hier nicht die Tote\pwindex{Schnitzler, Arthur 15.\,5.\,1862 Wien – 21.\,10.\,1931 ebd.@\textsc{Schnitzler, Arthur} (15.\,5.\,1862 Wien – 21.\,10.\,1931 ebd.), \emph{Schriftsteller, Mediziner}!Gefährtin. Schauspiel in einem Akt@\strich\emph{Die Gefährtin. Schauspiel in einem Akt}|pwv} ent{\pb}larven,{ }ſondern vielmehr die beiden Männer\pwindex{Schnitzler, Arthur 15.\,5.\,1862 Wien – 21.\,10.\,1931 ebd.@\textsc{Schnitzler, Arthur} (15.\,5.\,1862 Wien – 21.\,10.\,1931 ebd.), \emph{Schriftsteller, Mediziner}!Lebendige Stunden@\strich\emph{Lebendige Stunden}|pwv} erfaſſt haben, die Lebenden. Durch jene ganz
               unerwartete Wendung, mit welcher \textsc{Heinrich\pwindex{Schnitzler, Arthur 15.\,5.\,1862 Wien – 21.\,10.\,1931 ebd.@\textsc{Schnitzler, Arthur} (15.\,5.\,1862 Wien – 21.\,10.\,1931 ebd.), \emph{Schriftsteller, Mediziner}!Lebendige Stunden@\strich\emph{Lebendige Stunden}|pw}} die Laſt des Briefes gleichſam abwirft, indem er nachweist, daſs er nicht für
               ihn beſtimmt war, wird die Handlung ganz wieder den beiden Männern\pwindex{Schnitzler, Arthur 15.\,5.\,1862 Wien – 21.\,10.\,1931 ebd.@\textsc{Schnitzler, Arthur} (15.\,5.\,1862 Wien – 21.\,10.\,1931 ebd.), \emph{Schriftsteller, Mediziner}!Lebendige Stunden@\strich\emph{Lebendige Stunden}|pwv} gegeben und die Tote\pwindex{Schnitzler, Arthur 15.\,5.\,1862 Wien – 21.\,10.\,1931 ebd.@\textsc{Schnitzler, Arthur} (15.\,5.\,1862 Wien – 21.\,10.\,1931 ebd.), \emph{Schriftsteller, Mediziner}!Lebendige Stunden@\strich\emph{Lebendige Stunden}|pwv}
               hat keinen Theil mehr daran. Sie weiß nicht mehr davon. Sie iſt{ }ſchon überlebt,
               überwunden von dieſer einzigen Reflexbewegung des Lebens, die \textsc{Heinrich\pwindex{Schnitzler, Arthur 15.\,5.\,1862 Wien – 21.\,10.\,1931 ebd.@\textsc{Schnitzler, Arthur} (15.\,5.\,1862 Wien – 21.\,10.\,1931 ebd.), \emph{Schriftsteller, Mediziner}!Lebendige Stunden@\strich\emph{Lebendige Stunden}|pwv}} in{ }ſeiner Bedrängnis macht. Und aus dieſer Stimmung heraus begreift man den
               Schluſs in{ }ſeiner ganzen Einfachheit und Stärke. Sogar der Schauplatz ändert{ }ſich:
               der letzte Blick fällt auf den Garten!{\dots}\pend
           
\pstart
           Nun, verzeihen Sie, daſs ich mich von der{ }ſtarken Ergriffenheit verleiten ließ,{ }ſoviel über das kleine Stück\pwindex{Schnitzler, Arthur 15.\,5.\,1862 Wien – 21.\,10.\,1931 ebd.@\textsc{Schnitzler, Arthur} (15.\,5.\,1862 Wien – 21.\,10.\,1931 ebd.), \emph{Schriftsteller, Mediziner}!Lebendige Stunden@\strich\emph{Lebendige Stunden}|pwv} zu{ }ſchreiben; man thut gewöhnlich nicht gut daran,{ }ſich in der erſten Freude über die
               Urſachen eines Eindrucks klar werden zu wollen – {\pb}man kann dadurch doch nichts von der
               Freude, die man empfunden hat übertragen, auf den, der der mittelbare Schöpfer dieſer
               Freude war. Nun,{ }ſo kommt wenigstens das kleine beſcheidene \label{K_L04540-2v}\edtext{graue Buch\pwindex{Rilke, Rainer Maria 4.\,12.\,1875 Prag – 29.\,12.\,1926 Montreux@\textsc{Rilke, Rainer Maria} (4.\,12.\,1875 Prag – 29.\,12.\,1926 Montreux)!Letzten@\strich\emph{Die Letzten}|pwv}}{\lemma{\textnormal{\emph{graue Buch}}}\Cendnote{\textnormal{XXXX Auszeichnungsfehler: Dokument L04537 nicht gefunden.}}}\label{K_L04540-2}, das ich Ihnen in dieſen Tagen{ }ſchon immer{ }ſenden wollte,
               zugleich mit einer neuen Dankbarkeit zu Ihnen. Ich möchte, daſs es Ihnen etwas Freude
               bereitet! Außerdem erſcheint bei\textsc{Langen\orgindex{Albert Langen@Albert Langen|pw}} ein kleines zweiaktiges Drama »\textsc{Das tägliche Leben\pwindex{Rilke, Rainer Maria 4.\,12.\,1875 Prag – 29.\,12.\,1926 Montreux@\textsc{Rilke, Rainer Maria} (4.\,12.\,1875 Prag – 29.\,12.\,1926 Montreux)!tägliche Leben. Drama in zwei Akten@\strich\emph{Das tägliche Leben. Drama in zwei Akten}|pw}}«, das Sie auch bekommen{ }ſollen. Es{ }ſcheint indeſſen, daſs der Verlag\orgindex{Albert Langen@Albert Langen|pwv} die Buchausgabe \label{K_L04540-3v}\edtext{erſt am Tag der Erſtaufführung\eventindex{Residenz-Theater Berlin@\textbf{Residenz-Theater Berlin}!Uraufführung von Das tägliche Leben, 20.12.1901@Uraufführung von Das tägliche Leben, 20.12.1901|pwv} ausgeben}{\lemma{\textnormal{\emph{erst … ausgeben}}}\Cendnote{\textnormal{Das \emph{Börsenblatt für den deutschen Buchhandel}\pwindex{Börsenblatt für den Deutschen Buchhandel@\emph{Börsenblatt für den Deutschen Buchhandel}|pwk} meldet es
                        am 17. 12. 1901 verfügbar, drei Tage vor der Uraufführung\eventindex{Residenz-Theater Berlin@\textbf{Residenz-Theater Berlin}!Uraufführung von Das tägliche Leben, 20.12.1901@Uraufführung von Das tägliche Leben, 20.12.1901|pwkv}.}}}\label{K_L04540-3} will:{ }ſo müſſen Sie noch eine Weile warten. Das Stück\pwindex{Rilke, Rainer Maria 4.\,12.\,1875 Prag – 29.\,12.\,1926 Montreux@\textsc{Rilke, Rainer Maria} (4.\,12.\,1875 Prag – 29.\,12.\,1926 Montreux)!tägliche Leben. Drama in zwei Akten@\strich\emph{Das tägliche Leben. Drama in zwei Akten}|pw} ist für die litterar.
                        Abende des Reſidenztheaters\orgindex{Residenz-Theater Berlin@Residenz-Theater Berlin|pw} in \textsc{Berlin\oindex{Berlin@\textbf{Berlin}|pw}} und für das \textsc{Deutsche Schauspielhaus\orgindex{Deutsches Schauspielhaus in Hamburg@Deutsches Schauspielhaus in Hamburg|pw}} (Baron \textsc{Berger\pwindex{Berger, Alfred von 30.\,4.\,1853 Wien – 24.\,8.\,1912 ebd.@\textsc{Berger, Alfred von} (30.\,4.\,1853 Wien – 24.\,8.\,1912 ebd.)|pw}}) in \textsc{Hamburg\oindex{Hamburg@\textbf{Hamburg}|pw}} angenommen, aber über den Termin der Aufführung\eventindex{Residenz-Theater Berlin@\textbf{Residenz-Theater Berlin}!Uraufführung von Das tägliche Leben, 20.12.1901@Uraufführung von Das tägliche Leben, 20.12.1901|pwv}
               iſt mir noch nichts bekannt. Ich grüße Sie, lieber und verehrter Doctor \textsc{Schnitzler}, in viel aufrichtiger Herzlichkeit und Verehrung!\pend
           
\pstart
           Ihr:{\\[\baselineskip]}\spacefill\mbox{RainerMariaRilke}\pend
           \leftskip=0em{}\selectlanguage{ngerman}\endnumbering\briefempfaengerindex{Schnitzler, Arthur@\textsc{Schnitzler, Arthur}!zzzRilke, Rainer Maria@\emph{von Rainer Maria Rilke}!1901-12-031@{3. 12. 1901}|)be}\mylabel{L04540h}
\begin{anhang}
\end{anhang}\newcommand{\dateiname}{L04540}\newcommand{\titel}{Rainer Maria Rilke an Arthur Schnitzler, 3. 12. 1901}\newcommand{\editorInnen}{Herausgegeben von Jahnke, SelmaMüller, Martin Anton}\input{../tex-inputs/latex-leseansicht-abspann}
